\documentclass[12pt,a4paper]{ctexart}

% ============================================================
%                       宏包设置
% ============================================================

\usepackage{amsmath, amssymb, amsfonts}
\usepackage{geometry}
\usepackage{fancyhdr}
\usepackage{graphicx}
\usepackage{xcolor}
\usepackage{setspace}
\usepackage[most]{tcolorbox}
\usepackage{enumitem}

% ============================================================
%                 清新蓝绿色配色
% ============================================================

% 主色：清爽青绿色
\definecolor{SZURed}{HTML}{2F9E8F}

% 题目区域：浅薄荷绿
\definecolor{SZULightRed}{HTML}{EAF7F4}

% 答案区域：浅天空蓝
\definecolor{SZULightRose}{HTML}{EEF7FC}

% 边框：柔和蓝绿色
\definecolor{SZUSoftBorder}{HTML}{9CCFC8}

% 强调文字：沉稳蓝绿色
\definecolor{SZUTextRed}{HTML}{256B72}

% ============================================================
%                       页面设置
% ============================================================

\geometry{
    left=2.5cm,
    right=2.5cm,
    top=2.5cm,
    bottom=2.5cm
}

\setlength{\parindent}{0pt}
\setlength{\parskip}{0.6em}
\linespread{1.15}

% ============================================================
%                       课程信息
% ============================================================

\newcommand{\coursename}{概率论与数理统计}
\newcommand{\department}{深圳大学生物医学工程}
\newcommand{\homeworknumber}{作业 5}

% ============================================================
%                  学生信息 —— 请学生修改
% ============================================================

\newcommand{\studentname}{请填写姓名}
\newcommand{\studentid}{请填写学号}

% ============================================================
%                       页眉页脚
% ============================================================

\pagestyle{fancy}
\fancyhf{}

\lhead{\textcolor{SZUTextRed}{\coursename}}
\rhead{\textcolor{SZUTextRed}{\homeworknumber}}
\cfoot{\textcolor{SZUTextRed}{\thepage}}

\renewcommand{\headrulewidth}{0.4pt}
\renewcommand{\headrule}{
    \hbox to\headwidth{
        \color{SZUSoftBorder}\leaders\hrule height \headrulewidth\hfill
    }
}

\setlength{\headheight}{15pt}

% ============================================================
%                    题目编号与环境设置
% ============================================================

\newcounter{problem}

% -------------------------
% 题目环境
% -------------------------

\newenvironment{problem}[1][]
{
    \stepcounter{problem}
    \begin{tcolorbox}[
        title={题目 \theproblem\if&#1&\relax\else\quad[#1]\fi},
        colback=SZULightRed,
        colframe=SZURed,
        coltitle=white,
        fonttitle=\bfseries,
        boxrule=0.8pt,
        arc=2mm,
        left=3mm,
        right=3mm,
        top=2.5mm,
        bottom=2.5mm,
        breakable
    ]
}
{
    \end{tcolorbox}
}

% -------------------------
% 答案环境
% -------------------------

\newenvironment{answer}
{
    \begin{tcolorbox}[
        title={学生答案},
        colback=SZULightRose,
        colframe=SZUSoftBorder,
        coltitle=SZUTextRed,
        fonttitle=\bfseries,
        boxrule=0.7pt,
        arc=2mm,
        left=3mm,
        right=3mm,
        top=2.5mm,
        bottom=2.5mm,
        breakable
    ]
}
{
    \end{tcolorbox}
}

% ============================================================
%                         正文
% ============================================================

\begin{document}

% ============================================================
%                         标题
% ============================================================

\begin{center}

    {\LARGE\bfseries\textcolor{SZURed}{\coursename}}

    \vspace{0.4em}

    {\large \department}

    \vspace{1em}

    {\Large\bfseries\textcolor{SZUTextRed}{\homeworknumber}}

\end{center}

\vspace{1em}

% ============================================================
%                       学生信息
% ============================================================

\begin{center}

\begin{tabular}{rl@{\hspace{3cm}}rl}

\textcolor{SZUTextRed}{\bfseries 姓名：}
&
\studentname
&
\textcolor{SZUTextRed}{\bfseries 学号：}
&
\studentid

\end{tabular}

\end{center}

\vspace{0.7em}

{\color{SZUSoftBorder}\hrule}

\vspace{1.5em}



% ============================================================
% ============================================================
%
%                           题目 1
%
% ============================================================
% ============================================================

\begin{problem}[课本习题1]

在总体 $N(52,6.3^2)$ 中随机抽取一容量为 36 的样本，求样本均值 $\bar{X}$ 落在 50.8 到 53.8 之间的概率。
\end{problem}


% ============================================================
%
%              ↓↓↓↓↓↓ 学生答案区域 ↓↓↓↓↓↓
%
%          请只在 answer 环境中填写你的答案
%
% ============================================================

\begin{answer}

% ==========================================================
%                 请从下一行开始填写答案
% ==========================================================






% ==========================================================
%                 请在上一行结束填写答案
% ==========================================================

\end{answer}

% ============================================================
%
%              ↑↑↑↑↑↑ 学生答案区域 ↑↑↑↑↑↑
%
% ============================================================


\vspace{1em}



% ============================================================
% ============================================================
%
%                           题目 2
%
% ============================================================
% ============================================================

\begin{problem}[课本习题3]

求总体 $N(20,3)$ 的容量分别为 10,15 的两独立样本均值差的绝对值大于 0.3 的概率。
\end{problem}


% ============================================================
%
%              ↓↓↓↓↓↓ 学生答案区域 ↓↓↓↓↓↓
%
% ============================================================

\begin{answer}

% ==========================================================
%                 请从下一行开始填写答案
% ==========================================================





% ==========================================================
%                 请在上一行结束填写答案
% ==========================================================

\end{answer}

% ============================================================
%
%              ↑↑↑↑↑↑ 学生答案区域 ↑↑↑↑↑↑
%
% ============================================================


\vspace{1em}



% ============================================================
% ============================================================
%
%                           题目 3
%
% ============================================================
% ============================================================

\begin{problem}[课本习题8]

设总体 $X\sim N(\mu,\sigma^2)$，$X_1,X_2,\dots,X_{10}$ 是来自 $X$ 的样本。
\begin{enumerate}[label=(\arabic*)]
    \item 写出 $X_1,X_2,\dots,X_{10}$ 的联合概率密度。
    \item 写出 $\bar{X}$ 的概率密度。
\end{enumerate}
\end{problem}


% ============================================================
%
%              ↓↓↓↓↓↓ 学生答案区域 ↓↓↓↓↓↓
%
% ============================================================

\begin{answer}

% ==========================================================
%                 请从下一行开始填写答案
% ==========================================================





% ==========================================================
%                 请在上一行结束填写答案
% ==========================================================

\end{answer}

% ============================================================
%
%              ↑↑↑↑↑↑ 学生答案区域 ↑↑↑↑↑↑
%
% ============================================================


\vspace{1em}



% ============================================================
% ============================================================
%
%                           题目 4
%
% ============================================================
% ============================================================

\begin{problem}[课本习题10]

下面列出了30个美国NBA球员的体重（以磅计，1磅$=0.454\mathrm{kg}$）数据。这些数据是从美国NBA球队1990—1991赛季的花名册中抽样得到的。

\vspace{0.5em}
\centerline{225\quad 232\quad 232\quad 245\quad 235\quad 245\quad 270\quad 225\quad 240\quad 240}
\centerline{217\quad 195\quad 225\quad 185\quad 200\quad 220\quad 200\quad 210\quad 271\quad 240}
\centerline{220\quad 230\quad 215\quad 252\quad 225\quad 220\quad 206\quad 185\quad 227\quad 236}

\begin{enumerate}[label=(\arabic*)]
    \item 画出这些数据的频率直方图（提示：最大和最小观察值分别为 271 和 185，区间 $[184.5,271.5]$ 包含所有数据，将整个区间分为5等份，为计算方便，将区间调整为 $(179.5,279.5)$）。
    \item 作出这些数据的箱线图。
\end{enumerate}
\end{problem}


% ============================================================
%
%              ↓↓↓↓↓↓ 学生答案区域 ↓↓↓↓↓↓
%
% ============================================================

\begin{answer}

% ==========================================================
%                 请从下一行开始填写答案
% ==========================================================





% ==========================================================
%                 请在上一行结束填写答案
% ==========================================================

\end{answer}

% ============================================================
%
%              ↑↑↑↑↑↑ 学生答案区域 ↑↑↑↑↑↑
%
% ============================================================





\end{document}
```
